

\task{Selecting the freewheel diode for two different converters}

A manufacturer has two diode candidates for use as the freewheel element.
Following parameters are given:

\begin{table}[ht]
    \centering  % Zentriert die Tabelle
    \begin{tabular}{ll}
        \toprule
        Forward voltage $V_\mathrm{F}$ (at $I_\mathrm{F} = \SI{12}{\ampere}$) &  $\SI{1.05}{\volt}$\\ 
        Reverse charge $Q_\mathrm{rr}$ &  $\SI{4.2}{\micro\coulomb}$ \\ 
        Reverse recovery time $t_\mathrm{rr}$ &  $\SI{110}{\nano\second}$ \\ 
        Maximal reverse voltage  $V_\mathrm{R,\max}$ &  $\SI{1200}{\volt}$ \\
        \bottomrule
    \end{tabular}
    \caption{Device A: fast silicon p--n diode.}  % Beschriftung der Tabelle
    \label{table:ex01_deviceA}
\end{table}

and

\begin{table}[ht]
    \centering  % Zentriert die Tabelle
    \begin{tabular}{ll}
        \toprule
        Forward voltage $V_\mathrm{F}$ (at $I_\mathrm{F} = \SI{12}{\ampere}$) : &  $\SI{0.48}{\volt}$\\ 
        Reverse charge $Q_\mathrm{rr}$: &  negligible \\ 
        junction capacitance $C_\mathrm{j}$ &  $\SI{950}{\pico\farad}$ \\ 
        reverse leakage at room temperature $I_\mathrm{R}$: &  $\SI{2.5}{\milli\ampere}$ \\
        Maximal reverse voltage  $V_\mathrm{R,\max}$ &  $\SI{150}{\volt}$ \\
        \bottomrule
    \end{tabular}
    \caption{Device B: Schottky diode.}  % Beschriftung der Tabelle
    \label{table:ex01_deviceA}
\end{table}


Two applications are under consideration.
\textbf{Application I}
A $\SI{48}{\volt}$ dc-dc converter, output current $\SI{12}{\ampere}$, switching frequency $\SI{100}{\kilo\hertz}$.
\textbf{Application II}
A $\SI{700}{\volt}$ dc-link chopper leg, current $\SI{12}{\ampere}$, switching frequency $\SI{20}{\kilo\hertz}$.

\subtask{For each application, decide which diode should be selected and justify the choice physically.}

\begin{solutionblock}
    For Application I, choose Device B (Schottky). The bus is only $\SI{48}{\volt}$, so the $\SI{150}{\volt}$ blocking capability is acceptable. 
    At $\SI{100}{\kilo\hertz}$, reverse recovery is highly undesirable, and the Schottky has none essentially. 
    This matches the fact that Schottky diodes are majority-carrier devices, fast, low $V_F$, and preferred in low-voltage high-speed circuits.

    For Application II, choose Device A (fast silicon p--n diode). 
    The Schottky cannot be used because its reverse-voltage rating is only $\SI{150}{\volt}$, whereas the application needs $\SI{700}{\volt}$ blocking 
    with a margin. The poor high-voltage blocking of Schottky diodes and the trade-off between breakdown voltage and drift-region resistance.
\end{solutionblock} 

\subtask{Estimate the forward-conduction loss in each diode if it carries 12 A for 50\% duty:
    \begin{align*}
        P_{\mathrm{cond}} = V_\mathrm{F} \cdot I \cdot D
    \end{align*}}

\begin{solutionblock}
    With $I = \SI{12}{\ampere}$ and $D = 0.5$,

    For Device A:
    \begin{equation}
        P_{\mathrm{cond},A} = V_\mathrm{F} \cdot I \cdot D = \SI{1.05}{\volt} \cdot \SI{12}{\ampere} \cdot 0.5 = \SI{6.3}{\watt}.
    \end{equation}
    For Device B:
    \begin{equation}
        P_{\mathrm{cond},B} = = V_\mathrm{F} \cdot I \cdot D \SI{0.48}{\volt} \cdot \SI{12}{\ampere} \cdot 0.5 = \SI{2.88}{\watt}
    \end{equation}
    So the Schottky saves:
    \begin{equation}
        P_{\mathrm{cond},A} - P_{\mathrm{cond},B} = \SI{6.3}{\watt} - \SI{2.88}{\watt} = \SI{3.42}{\watt}.
    \end{equation}
\end{solutionblock} 


\subtask{For Device A, estimate the reverse-recovery loss in Application II using
    \begin{align*}
        E_\mathrm{rr} &\approx Q_\mathrm{rr} V_\mathrm{R} \\
        P_\mathrm{rr} &= E_\mathrm{rr} f_\mathrm{s}
    \end{align*}
}


\begin{solutionblock}
    Reverse-recovery loss of Device A in Application II is calculated by
    \begin{equation}
        E_\mathrm{rr} \approx Q_\mathrm{rr} V_\mathrm{R}  = \SI{4.2}{\micro\coulomb} \cdot \SI{700}{\volt} = \SI{2.94}{\milli\joule}.
    \end{equation}
    \begin{equation}
        P_\mathrm{rr} = E_\mathrm{rr} f_\mathrm{s} = \SI{2.94}{\milli\joule} \cdot \SI{20}{\kilo\hertz} = \SI{58.8}{\watt}.
    \end{equation}

    This is much larger than the conduction loss.
\end{solutionblock} 

\subtask{Explain why Device B may still be rejected for Application II even though it has lower forward loss and no reverse recovery.}

\begin{solutionblock}
    Device B is still rejected in Application II, because device selection is not determined only by low $V_F$ or fast switching. 
    The diode must survive the blocking state. Device B fundamentally cannot support the required reverse voltage. 
    This is the central trade-off in power-device design: majority-carrier speed helps switching, 
    but high-voltage blocking usually requires a different structure.
\end{solutionblock} 

\subtask{State one situation in which Device A is the only realistic option, and one situation in which Device B is clearly preferable.}

\begin{solutionblock}
    The final design judgment bases on following points:
    \begin{itemize}
        \item Device A only realistic: high-voltage freewheel path, e.g. $\SI{700}{\volt}$ -- $\SI{1200}{\volt}$ leg.
        \item Device B clearly preferable: low-voltage, high-frequency output rectifier.
    \end{itemize} 
\end{solutionblock} 

\task{The drift region as the real blocking element of a power diode}

Consider a vertical one-sided silicon power diode of type $p^+ - n^- - n^+$. The $n^-$ drift region is designed to support reverse voltage.


\begin{table}[ht]
    \centering  % Zentriert die Tabelle
    \begin{tabular}{ll}
        \toprule
        Donator density $N_\mathrm{D}$          &  $\SI{10^{14}}{\centi\meter^{-3}}$ \\ 
        Drift-region thickness $t_\mathrm{d}$   &  $\SI{80}{\micro\meter}$ \\ 
        Reluction time $\varepsilon_\mathrm{s}$ &  $\SI{1.04 \cdot 10^{-12}}{\farad\centi\meter^{-1}}$ \\ 
        Elementary charge $q$                   &  $\SI{1.6\cdot 10^{-19}}{\coulomb}$ \\
        Electron mobility $\mu_n$ in the drift region & $\SI{1000}{\centi\meter^2\per\volt\per\second}$ \\
        \multicolumn{2}{l}{built-in voltage may be neglected compared to large reverse voltage} \\
        \bottomrule
    \end{tabular}
    \caption{Device A: fast silicon p--n diode.}  % Beschriftung der Tabelle
    \label{table:ex01_deviceA}
\end{table}


\subtask{Under reverse bias, explain qualitatively why the depletion region extends mainly into the $n^-$ side rather than the $p^+$ side.}

\begin{solutionblock}
    The depletion region extends mainly into the $n^-$ side, because the $p^+$ side is very heavily doped compared with the $n^-$ drift region. 
    Charge neutrality requires equal uncovered charge on both sides, so the lightly doped side must widen much more. 
    This is the same electrostatic principle used in ordinary abrupt junctions, but it becomes the key design mechanism in power devices.
\end{solutionblock} 

\subtask{Estimate the depletion width at $V_\mathrm{R} = \SI{600}{\volt}$ using
    \begin{equation*}
        W = \sqrt{\frac{2 \varepsilon_\mathrm{s} V_\mathrm{R}} {q N_\mathrm{D}}}
    \end{equation*}
}

\begin{solutionblock}
    The depletion width $W$ is obtained by
        \begin{align}   
            W &= \sqrt{\frac{2 \varepsilon_\mathrm{s} V_\mathrm{R}} {q N_\mathrm{D}}}
              = \sqrt{\frac{2 \cdot \SI{1.04 \cdot 10^{-12}}{\farad\per\centi\meter} \cdot \SI{600}{\volt}} {\SI{1.6\cdot 10^{-19}}{\coulomb} \cdot \SI{10^{14}}{1\per\centi\meter^{3}}}} \\
              &= \sqrt{\frac{\SI{1.248 \cdot 10^{-9}}{\volt\farad\per\centi\meter}} {\SI{3.2 \cdot 10^{-5}}{\coulomb\per\centi\meter^{3}}}} = \SI{62.4}{\micro\meter}
        \end{align}             
\end{solutionblock} 

\subtask{Decide whether the diode is still non-punch-through at $\SI{600}{\volt}$ by comparing $W$ and $t_\mathrm{d}$.}

\begin{solutionblock}
    If the drift-region thickness $t_\mathrm{d}$ is greater than depletion width $W$, then a punch-through is prevented.
    This is the case: $\SI{80}{\micro\meter} > \SI{62.4}{\micro\meter}$.
\end{solutionblock} 


\subtask{Estimate the peak electric field using
    \begin{equation*}
        E_\mathrm{\max} = \frac{q N_\mathrm{D} W}{\varepsilon_\mathrm{s}}
    \end{equation*}
}

\begin{solutionblock}
    The peak electric field is estimated by
        \begin{equation}   
            E_\mathrm{\max} = \frac{q N_\mathrm{D}}{\varepsilon_\mathrm{s}}
             = \frac{\SI{1.6\cdot 10^{-19}}{\coulomb} \cdot \SI{10^{14}}{\centi\meter^{-3}}}{\SI{1.04 \cdot 10^{-12}}{\farad\centi\per\meter}}
             = \SI{1.92\cdot 10^{5}}{\volt\per\centi\meter}
        \end{equation}             
\end{solutionblock} 


\subtask{Now double the drift doping to $\SI{4 \cdot 10^{14}}{1\per\centi\meter}^{3}$, keeping thickness unchanged. 
         Without recalculating everything from scratch, state how $W$, $E_{\max}$, blocking capability, and drift-region resistance change.
}

\begin{solutionblock}
    The depletion width $W$ is proportional to the inverted square root of the donator density:
        \begin{equation}   
            W \propto \frac{1}{\sqrt{N_\mathrm{D}}}
        \end{equation}     
    so that $W$ decreases by $\sqrt{2}$. On the other hand The peak electric field increases by $\sqrt{2}$ according
        \begin{equation}   
            E_\mathrm{\max} \propto \sqrt{N_\mathrm{D}}
        \end{equation}    

    Consequences:
    \begin{itemize}
        \item depletion width decreases
        \item peak field increases
        \item blocking capability worsens
        \item drift resistance decreases
    \end{itemize}
\end{solutionblock} 

\subtask{ Use
        \begin{equation*}   
            R_{\mathrm{drift}} \propto \frac{t_\mathrm{d}}{q \mu_\mathrm{n} N_\mathrm{D}}
        \end{equation*}  
    to explain the statement: ``High $V_\mathrm{BD}$ $\Rightarrow$ low doping $\Rightarrow$ high $R_\mathrm{on}$.''
}

\begin{solutionblock}
    To support a large reverse voltage without reaching the critical field too early, the drift region must be lightly doped. 
    But lower $N_D$ means fewer carriers and thus larger resistivity:
        \begin{equation*}   
            \rho \propto \frac{1}{q \mu_\mathrm{n} N_\mathrm{D}}
        \end{equation*}  
    Therefore, the same region that helps blocking makes on-state conduction worse. 
    That is the structural origin of the power-device trade-off seen.
\end{solutionblock} 

